\begin{activity}{Stefan Problems and Cooling of Magma Bodies}

\section*{Purpose}
To apply transient heat flow concepts to the crystallization and cooling of magma bodies within the crust.

\section*{Physical Problem}
A magma body emplaced into cooler host rock cools conductively while releasing
latent heat at the solidification front.

\section*{Key Feature}
The solid--liquid boundary migrates with time, making this a Stefan problem.

\section*{Tasks}
\begin{enumerate}
  \item Explain how latent heat affects the cooling rate of a magma body.
  \item Compare expected cooling timescales for a narrow dike and a thick sill.
  \item Identify geological situations in which analytic Stefan solutions are
  likely to fail.
\end{enumerate}

\section*{Geological Extensions}
\begin{itemize}
  \item Magma convection.
  \item Episodic recharge of magma.
  \item Crystal-rich mush zones.
\end{itemize}

\section*{Key Insight}
Geometry and diffusivity often dominate cooling timescales more strongly than
radiogenic heat production.

\end{activity}